\cotefile{156-7}
\scanpage{87}
\mspage{86}

% slanted
\margin{\uline{26.6} Finalement, cet axiome m'apparaît un peu artificiel, il n'est pas satisfait par les ``gros objets'' \unsure{envisagés} : un espace top. $V$, \unsure{pour les} \unsure{variétés}, \unsure{A|à} moins de dire que c'est en gros \unsure{$\perp_C$|$\perp$} qui est ``artificiel'' ! Mais ce n'est pas \unsure{satisfait|satisfaisant} \unsure{non plus pour l'adhérence|non plus pour l'axiome}, \unsure{limitation par manque d'objets} topologiques\ldots}

\unsure{Je viens|Je vais} d'introduire un axiome \unsure{vraiment} crucial :

\uline{At supp 1} : Soient $X,Y\in\unsure{M|\mathcal M}$, supposons que
\[
Y\in X^\circ\overset{\mathrm{def}}{\equiv}\operatorname{supp}X\cap\operatorname{cosupp}\partial X,\quad\text{i.e. supposons}
\]
\[
\begin{aligned}
\text{(i)}\quad &Y\parallel\partial X &&\text{i.e.}\quad\forall Y'<X,\quad Y'\parallel Y\\
\text{(ii)}\quad &Y\in\operatorname{supp}_{M}X &&\text{i.e.}\quad\forall Z\in M,\ Z\parallel X\Rightarrow Z\parallel Y.
\end{aligned}
\]
Alors $Y\ll X$ \deleted{$\unsure{\Longleftrightarrow}$}

\deleted{\uline{Corollaire} Soient $F$ une figure, $G$ une sous-figure. Alors pour la \deleted{figure} $X(\in M)$, les conditions suivantes sont équivalentes}
\[
\deleted{\text{(i)}\quad H\ll F,}
\]

\uline{Corollaire}
\[
X^\circ=\left\{Y\in\operatorname{Omb}(X)\ \middle|\ \begin{aligned}Y_{\partial X}&=\varnothing,\ \text{i.e.}\\ \widetilde Y\cap\operatorname{Omb}(\partial X)&=\varnothing\end{aligned}\right\}
\]

En effet, si $Y\deleted{\unsure{\in}}\ll X$, \deleted{\unsure{la relation} $Y_{\partial X}=\varnothing$,} on sait par At D que
\[
\text{\unsure{car|or|on}}\quad Y_{\partial X}=\varnothing\Longleftrightarrow Y\parallel\partial X
\]
(on a $\Leftarrow$ sous hyp. At D) donc l'\unsure{ens} du second membre \deleted{\unsure{est}} contient $X^\circ$ par Absorpt. L'inclusion inverse, compte tenu de (i), est triviale ($\mathrm{NB}\ Y\in\operatorname{Omb}(X)\Rightarrow Y\in\operatorname{supp}_{M}X$).
\scanpage{88}
\mspage{87}

\uline{Remarque} On peut songer à renforcer Absorp 1 ainsi :

\uline{Absorp 1} (spec) Soient $X,Y,\deleted{\unsure{Y'}}\in\unsure{M|\mathcal M}$, avec $\deleted{Y'}\inserted{\unsure{\widetilde L|L}}\leq Y$.

Supposons
\[
\begin{aligned}
\text{(i)}\quad &\operatorname{supp}Y\subset X\\
\text{(ii)}\quad &\deleted{\unsure{Y'}}L\ll\partial X\\
\text{(iii)}\quad &\text{``}Y\setminus\deleted{Y'}L\text{''}\parallel\partial X,\ \text{de façon précise,}\quad\forall Z\in\operatorname{Omb}(Y),\\
&Z\deleted{\unsure{\cap}}_{\deleted{\unsure{Y'}}L}=\varnothing\Rightarrow Z\parallel\partial X
\end{aligned}
\]
Alors $Y\ll X$ (et $L=Y_{\partial X}$)

Ceci suggère qu'il doit exister une définition récurrente de la relation $\ll$ dans $M$, au moyen du $\leq$ et de la relation $\parallel_M$ -- mais cette dernière relation non seulement pour \unsure{deux objets} de $M$ : $X,Y$, mais aussi pour des différences virtuelles
\[
Y\setminus L\parallel\deleted{\unsure{\S|\partial}}Z
\]
\margin{$Y,Z\in M$, \deleted{$\unsure{L\in M},L\leq Y$} \inserted{$L$ partie fermée de $\widetilde Y$}}
voire même
\[
Y\setminus L\parallel Z\setminus K\qquad -
\]
cette relation intervient tacitement dans l'énoncé de At C (spec), mais circonscrite à coups de $\deleted{\unsure{Z}}S\ll Y\setminus L,\ T\ll Z\setminus K$ tels que\ldots
\scanpage{89}
\mspage{88}

Faisons abstraction pour le moment des axiomes de support posés, et de la définition proposée de $X^\circ$. Posons

% The block from Définition through the last display is crossed by long strokes.
\deleted{\uline{Définition} (i) \deleted{Soit} $X\deleted{\unsure{Y}}\in\unsure{M|\mathcal M}$. Posons}
\[
\deleted{\deleted{\unsure{X^\circ\parallel Y^\circ}}\quad\operatorname{Omb}^{\unsure{\circ\circ|\infty}}(X)=\{Y\in\operatorname{Omb}(X)\mid Y\parallel\partial X\}}
\]
\deleted{Evidemment}
\[
\deleted{\operatorname{Omb}^{\circ\circ}(X)\subset\operatorname{Omb}^{\circ}(X)=\operatorname{Omb}(X)\setminus\operatorname{Omb}(\partial X)}
\]

\deleted{(ii) Posons}
\[
\deleted{X^\circ\parallel Y^\circ\overset{\mathrm{def}}{\Longleftrightarrow}\operatorname{Omb}^{\circ\circ}(X)\cap\operatorname{Omb}^{\circ\circ}(Y)=\varnothing}
\]
\deleted{i.e. il n'existe pas de $Z\in M$ avec}
\[
\deleted{Z\ll X,Y,\quad Z|\partial X=\varnothing,\quad Z|\partial Y=\varnothing}
\]
\deleted{ou encore, $\deleted{\unsure{\forall}}$}
\[
\deleted{Z\ll X,Y\Rightarrow\exists Z'\in\widetilde Z\ \text{avec}\ Z'\unsure{<|\ll}\deleted{\unsure{\partial}}X\ \text{ou}\ Z'<Y.}
\]

\deleted{On aimerait avoir, pour $X,Y\in M$}
\[
\deleted{X\parallel Y\Longleftrightarrow\forall X'\in\widetilde X,Y'\in\widetilde Y,\quad X'^\circ\unsure{\nparallel|\parallel}Y'^\circ}
\]
\deleted{On a toujours $\Rightarrow$, car $X\parallel Y$ implique $\nexists Z$ avec $Z\ll X,Y$ donc a fortiori $\nexists Z$ avec $Z\ll X'^\circ,Y'^\circ$ (si $X'\in\widetilde X,Y'\in\widetilde Y$). L'implication plus difficile est $\Leftarrow$. Elle est conséquence de At L2 et At L3, car}
\[
\deleted{X\parallel Y\Longleftrightarrow\forall x\in\operatorname{omb}(X);\ \forall y\in\operatorname{omb}(Y),\quad x\parallel y}
\]
\scanpage{90}
\mspage{89}

\uline{Définition} Soit $X\in\unsure{M|\mathcal M}$. Posons
\[
X^\circ=\left\{Y\in\operatorname{Omb}(X)\ \middle|\ \begin{aligned}
&Y\parallel\partial X\quad\text{i.e. }Y_{\partial X}=\varnothing\quad\text{i.e.}\\
&\deleted{\unsure{\exists Z}}\widetilde Y\cap\operatorname{Omb}(\partial X)=\varnothing\\
&\text{i.e. }\unsure{\exists|\nexists}Y'\leq Y,\ Y'\ll X'<X
\end{aligned}\right\}
\]
\margin{ou $\operatorname{Omb}^{\unsure{\circ\circ|\infty}}(X)\subset\operatorname{Omb}^{\circ}(X)$}
% slanted
\margin{partie fermée de $M,\ll$}

Donc pour $X,Y\in M$, la relation
\[
X^\circ\parallel Y^\circ
\]
\unsure{est} définie, elle signifie
\[
\forall X'\in X^\circ,Y'\in Y^\circ,\quad\text{on a}\quad X'\parallel Y'.
\]
\margin{Attention, \uline{\unsure{je ne|ne me}} \uline{dis pas} que $X^\circ\cap Y^\circ=\varnothing\Rightarrow X^\circ\parallel Y^\circ$, cela \unsure{impliquerait} pour $x,y\in\unsure{L|\ell}$ $x\ne y\Rightarrow x\parallel y$ !}
% slanted
\margin{\uline{NB} $Y\in X^\circ$ signifie \unsure{aussi que} $Y\ll X$, et l'application $\unsure{\varphi|\phi}:\widetilde Y\to\widetilde X$ est constante de valeur le \deleted{\illegible{1}} plus grand élément de $X$. Au lieu de $Y\in X^\circ$, on écrit aussi $Y\in\operatorname{Omb}^{\circ\circ}(X)$, ou $Y\overset{\circ\circ}{\ll}X$ (\unsure{raffinement} \uline{intérieur} de $X$, plus fort que \unsure{raffinement} \uline{immédiat} $Y\overset{\circ}{\ll}X$ i.e. $Y\in\operatorname{Omb}^{\circ}(X)$).}

On voudrait prouver ceci

\inserted{(OK si on suppose At L2, L3)}

\uline{Lemme 1} (?) Soient $X,Y\in M$. Alors
\[
X\parallel Y\Longleftrightarrow\forall X'\in\widetilde X,Y'\in\widetilde Y,\quad\text{on a}\quad X'^\circ\parallel Y'^\circ.
\]

\uline{Dém.} $\Rightarrow$ est clair, car $X'^\circ\subset\operatorname{Omb}(X')\subset\operatorname{Omb}(X)$ et de $=$ $Y'^\circ\unsure{\subset|\in}\operatorname{Omb}(Y')\subset\operatorname{Omb}(Y)$, et on applique At D$_0$.

L'implication $\Leftarrow$ est moins claire, \deleted{Elle} l'est \deleted{\unsure{mais est loin alors à prouver ceci|mais est loin alors d'être claire}} car on doit \unsure{en admettre} L2, L3.
% slanted
\margin{\uline{Cor} $X\parallel Y\Longleftrightarrow\forall Y'\in\widetilde Y,\ X\parallel Y'^\circ$}

\uline{Lemme 2} \inserted{(At \textcircled{\small L1}, L2, L3)} \deleted{Supposons} Soit $Y\ll X$. Alors
\[
\text{a)}\quadY^\circ\overset{\circ}{\ll}X\quad\text{sss}\quad Y^\circ\subset X^\circ.\quad\text{On a}
\]
\[
\text{b)}\quad Y\overset{\circ\circ}{\ll}X\quad\inserted{\text{i.e. }Y\in X^\circ}\quad\text{sss}\quad\forall Y'\in\widetilde Y,\quad Y'^\circ\subset X^\circ.
\]

\uline{Dém.} a) Supposons $Y\overset{\circ}{\ll}X$, prouvons $Y^\circ\subset X^\circ$. Soit $Y'\in Y^\circ$, i.e. $Y'\ll Y$, $Y'\parallel\partial Y$, donc $Y'\ll X$, il faut prouver $Y'\parallel\partial X$ ou encore $Y'_{\partial X}=\varnothing$
\scanpage{91}
\mspage{90}

\diagram{\begin{tikzcd}[ampersand replacement=\&,row sep=small,column sep=small]
Y'\arrow[r,phantom,"\ll"{description}]\&Y\arrow[r,phantom,"\ll"{description}]\arrow[d,phantom,"\ll"{description,rotate=90}]\&X\arrow[d,phantom,"\ll"{description,rotate=90}]\\
\&Y_{\partial X}\arrow[r,phantom,"\ll"{description}]\arrow[d,phantom,"\ll"{description,rotate=90}]\&\partial X\\
\&\partial Y\&
\end{tikzcd}}

Or
\[
\begin{aligned}
&\deleted{\unsure{Y'_{\partial X}}}Y'|\partial X\quad(\text{induction dans }X)=\\
&\deleted{Y'_{\partial X}}Y'|(Y_{\partial X}\deleted{\unsure{\vert\partial X}})\quad(\text{induction dans }Y)
\end{aligned}
\]
\deleted{\unsure{grâce à}} \deleted{\unsure{l'hypothèse} :}

Comme $Y_{\partial X}\leq\partial Y$ (par hyp. $Y\overset{\circ}{\ll}X$), on a $Y'|Y_{\partial X}\leq Y'|\partial Y$, \deleted{\unsure{par}} or $Y'|\partial Y\unsure{=|\simeq}Y'_{\partial Y}=\varnothing$ par hyp. $Y'\in Y^\circ$, d'où $Y'|Y_{\partial X}=\varnothing$ qed.

Supposons $Y^\circ\subset X^\circ$ prouvons $Y\overset{\circ}{\ll}X$, i.e. qu'on n'a \uline{pas} $Y\ll\partial X$. Il faut utiliser ici
\[
\boxed{Y^\circ\ne\varnothing}
\]
(ce qui résulte de l'axiome de divisibilité L1, compte tenu de L3, et est de nature analogue à L1\ldots) Soit $Z\in Y^\circ$, donc $Z\in X^\circ$ par $Y^\circ\subset X^\circ$, donc $Z\parallel\partial X$. Si on avait $Y\ll\partial X$ d'où $Z\ll\partial X$, on aurait $Z\parallel Z$ par At D$_0$, absurde.
% slanted
\margin{\uline{NB} Moyennant des hyp. de \unsure{dimension finie}, la condition ``$Y^\circ\ne\varnothing\quad\forall Y\in\unsure{M|\mathcal M}$'' \uline{\unsure{équivaut}} à L1\ldots}

b) Supposons $Y\in X^\circ$, prouvons $\forall Y'\in\widetilde Y,Y'^\circ\subset X^\circ$. En \deleted{\unsure{effet}} fait, non seulement $Y'^\circ\subset\operatorname{Omb}(Y')$ est $\subset X^\circ$, mais $\operatorname{Omb}(Y')\subset\operatorname{Omb}(Y)$ l'est, \deleted{\unsure{i.e.}} et en fait \deleted{On} On a
\[
Y\in X^\circ\Rightarrow\operatorname{Omb}(Y)\subset X^\circ
\]
i.e. $X^\circ$ est une partie de $\deleted{\unsure{\operatorname{Omb}(X)}}\inserted{M}$ fermée pour $\ll$, comme on a \unsure{noté} déjà ant. Donc $\Rightarrow$ de b) est triv. Prouvons $\Leftarrow$, i.e.
\scanpage{92}
\mspage{91}

supposons $\displaystyle\bigcup_{Y'\in\widetilde Y}Y'^\circ\subset X^\circ$, prouvons $Y\in X^\circ$. i.e. $Y_{\partial X}=\varnothing$, \deleted{Soit $Z$} i.e. $\nexists Z\in\widetilde Y$, avec $Z\ll\partial X$. En effet, pour un tel $Z$, \deleted{on} \deleted{la relation} On aurait par hyp. $Z^\circ\subset X^\circ$, donc $Z\overset{\circ}{\ll}X$ (par a)), \uline{i.e.} $Z\nless\partial X$.

On gagne.

\uline{NB} Soient $X,Y\in\unsure{M|\mathcal M}$, $X\unsure{\ngeq|\gtrless}Y$ \inserted{$X\ne Y$}. Alors
\[
X^\circ\cap Y^\circ=\varnothing
\]
% slanted
\margin{Mieux, $X^\circ\parallel Y^\circ$ cf p. 95}

En effet, soit $L=X\cap Y$. Si on avait $Z\in X^\circ\cap Y^\circ\subset\operatorname{Omb}(X)\cap\operatorname{Omb}(Y)=\operatorname{Omb}(L)$, on aurait \deleted{donc} $Z\ll L$. Or on ne peut avoir à la fois $L=X$, $L=Y$ (puisque $X\ne Y$), disons $L\ne X$ donc $L\subset\partial X$. Donc $Z\ll\partial X$, et comme $Z\parallel\partial X$ $Z\parallel Z$, absurde !

\deleted{On a envie de} Soit maintenant $\Phi\subset M$ une \deleted{\unsure{classe}} partie fermée de $M$, fermée \unsure{de plus} $\leq$, formée des \deleted{objets} \inserted{multistrates} donc : donc à comp. libres. \unsure{sss} $\Phi\unsure{\simeq|=}\widetilde F$ (\unsure{quitte à} \unsure{agrandir} $F$, \unsure{interprété} comme une des \unsure{parties} de $M$ -- mais \unsure{la|les|le},
\scanpage{93}
\mspage{92}

développements actuels ne dépendent de $\unsure{I|F}$ que par l'intermédiaire des Strates $\unsure{M|\mathcal M}$, $\leq$, $\ll$, $\unsure{\ngeq|\gtrless}$.

\deleted{Il sera commode de \unsure{désidentifier}, et d'introduire un ens. d'indices ordonné $I$, pour \unsure{indexer} $\widetilde F$}

Pour toute partie finie $A$ de $\widetilde F=\Phi$, posons
\[
\begin{aligned}
S_A&=\operatorname{Supp}_{M}\left(\bigcup_{X\in A}X^\circ\right)\quad\in\overline\Sigma_M\\
&=\underset{X\in A}{\operatorname{Sup}_{\Sigma_M}}\operatorname{Supp}(X^\circ).
\end{aligned}
\]
\margin{$\deleted{\unsure{\Rightarrow}}$}

(\uline{NB} Si on admet At supp. 1, on a $\operatorname{Supp}X^\circ=X^\circ$, i.e. $X^\circ$ est déjà une partie-support de $M$.)

Par définition et par associativité des Sup, on aura donc, si $A=\bigcup_i A_i$
\[
S_A=\operatorname{Sup}_{\deleted{\unsure{i}}\Sigma_M}S_{A_i}.
\]
On a ainsi une application
\diagram{\begin{tikzcd}[ampersand replacement=\&,column sep=large]
\unsure{\mathscr P|\mathcal P}(\widetilde F\unsure{]|)}\arrow[r]\&\Sigma_M\ ;\&A\arrow[r,mapsto]\&S_A
\end{tikzcd}}
de l'ens. ordonné des parties de $\widetilde F$ dans l'ens \deleted{des} ordonné des supports pour $M$, application évidemment croissante.
